\begin{userprompt}
in hyperband tuner, how are number of trials decided ? I have set max\_epoches to 15 but my trials have been past 30 now
\end{userprompt}

\begin{botresponse}
In Hyperband tuning, the number of trials isn't directly set but is determined by how the algorithm allocates resources across different configurations.

Hyperband works by using a bracket-based approach where it starts with many configurations with few resources (epochs), then progressively filters out poorly performing ones while giving more resources to promising configurations.

The total number of trials depends on:

With max\_epochs = 15, you're seeing more than 30 trials because Hyperband generates multiple configurations in early brackets. For example, in the first bracket it might try 27 configurations with 1 epoch each, then keep the top 9 for 3 epochs, then the top 3 for 9 epochs, etc.

The formula for total trials is approximately:


Is this a Keras Hyperband implementation, or are you using another framework? That might help me give you more specific information about your particular setup and why you're seeing 30+ trials.

\begin{enumerate}
\item The max\_epochs parameter (your budget per configuration)
\item The reduction factor eta (typically 3)
\item The number of brackets s\_max + 1 (calculated from max\_epochs and eta)
\end{enumerate}
\end{botresponse}

\begin{userprompt}
no it is a keras implementation

this is the code

\# helper function to create the Tuners, run experiment and retrain the model as visualize f2

def create\_tuners(input\_shape, num\_classes, max\_trials=3):

\# Create directories for storing results

    os.makedirs(f'random\_tuner', exist\_ok=True)
    os.makedirs(f'hyperband\_tuner', exist\_ok=True)
    os.makedirs(f'bayes\_tuner', exist\_ok=True)

\# Define the build\_model function with preset input\_shape and num\_classes
    def hp\_build\_model(hp):
        return build\_model(hp, input\_shape, num\_classes)

\# Create the tuners

    tuner\_rs = RandomSearch(
        hp\_build\_model,
        objective=keras\_tuner.Objective("val\_f2\_macro", direction="max"),
        max\_trials=max\_trials,
        executions\_per\_trial=1,
        directory=f'random\_tuner',
        project\_name='ordinal\_severity'
    )

tuner\_hb = Hyperband(
        hp\_build\_model,
        objective=keras\_tuner.Objective("val\_f2\_macro", direction="max"),
        max\_epochs=max\_trials,
        factor=3,
        directory=f'hyperband\_tuner',
        project\_name='ordinal\_severity'
    )

tuner\_bo = BayesianOptimization(
        hp\_build\_model,
        objective=keras\_tuner.Objective("val\_f2\_macro", direction="max"),
        max\_trials=max\_trials,
        num\_initial\_points=2,
        directory=f'bayes\_tuner',
        project\_name='ordinal\_severity'
    )

return tuner\_rs, tuner\_hb, tuner\_bo

def run\_tuner(tuner, tuner\_name, epochs=3):
    print(f"\textbackslash\{\}nStarting \{tuner\_name\} search...")

\# Create callback
    f2\_callback = F2Callback(
        validation\_data=(X\_val\_processed, y\_val\_cumulative),
        raw\_targets=y\_val,
        num\_classes=num\_classes
    )

\# Search
    tuner.search(
        X\_train\_processed, y\_train\_cumulative,
        epochs=epochs,
        validation\_data=(X\_val\_processed, y\_val\_cumulative),
        callbacks=[f2\_callback]
    )

\# Get the best hyperparameters
    best\_hps = tuner.get\_best\_hyperparameters(num\_trials=1)[0]
    print(f"\textbackslash\{\}nBest Hyperparameters from \{tuner\_name\}:")

return best\_hps

def model\_retrain\_best(best\_hps):

model = build\_model(best\_hps, X\_train\_processed.shape[1], num\_classes)

f2\_callback = F2Callback(
        validation\_data=(X\_val\_processed, y\_val\_cumulative),
        raw\_targets=y\_val,
        num\_classes=num\_classes
    )

history = model.fit(
        X\_train\_processed, y\_train\_cumulative,
        epochs=5, \# the model f2 score does not improve drastically, there for epoches are capped at just 5
        validation\_data=(X\_val\_processed, y\_val\_cumulative),
        callbacks=[f2\_callback],
        batch\_size=1024,
        verbose=0
    )

return model, history

def visualize\_f2(best\_hps):
    trial\_result = \{
    'learning\_rate': best\_hps['learning\_rate'],
    'num\_layers': best\_hps['num\_layers'],
    'f2\_scores': f2\_callback.f2\_scores,
    'trial\_id': 'best\_trial'  \# or any label you want
    \}

df = pd.DataFrame(\{
        'epoch': list(range(1, len(trial\_result['f2\_scores']) + 1)),
        'val\_f2\_macro': trial\_result['f2\_scores'],
        'trial\_id': trial\_result['trial\_id']
    \})

sns.lineplot(data=df, x='epoch', y='val\_f2\_macro', hue='trial\_id')
    plt.title("F2 Score over Epochs for Best Trial")
    plt.show()

tuner\_rs, tuner\_hb, tuner\_bo = create\_tuners(X\_train\_processed.shape[1], num\_classes, max\_trials=15)

for each model there are only 5 training epoches
\end{userprompt}